\chapter{Experiments}
\label{ch:experiments}

Every number below is produced by \code{make pipeline} from a fixed seed and rendered into this
chapter by \code{scripts/results\_to\_tex.py}. Change an equation, change the module, rerun, and
the table changes.

\section{Protocol}
A single synthetic twin (Chapter~\ref{ch:domain}) is generated. Episodes of $T$ steps are
simulated; a fraction are nominal, the rest carry one fault drawn from the catalogue of
Chapter~\ref{ch:propagation} with a random origin of the matching type, a random onset in
$[T/4,T/2)$, and a random severity and ramp. Episodes are split 60/15/25 into train, calibration
and test. The model of Chapter~\ref{ch:symmetry} is trained with the loss \eqref{eq:loss}.
Every detector, the GNN field detector \eqref{eq:detection}, the service-layer $z$-score trigger,
and the per-node rolling $z$-score, is calibrated on the \emph{same} nominal calibration episodes
with the same margin, so each has zero false alarms on calibration data by construction. The
static threshold is fixed at $2.5$ standardised units.

Metrics (Chapter~\ref{ch:closed} and below): a detector's delay is the first firing at or after
onset; a firing before onset is counted as a false alarm. Lead time is reported only over episodes
where the baseline fired, and the detection-rate gap is reported separately.
\begin{align}
\text{lead}&=t_{\text{baseline}}-t_{\text{gnn}} \label{eq:leadtime}\\
\text{hit@}k&=\ind{v_0\in\text{top-}k\text{ of the ranking at }t_{\text{gnn}}+\delta} \label{eq:hitk}\\
\text{triage@}k&=1-\frac{k}{|\{u:\ d_u(t_{\text{alarm}}+\delta)>0.5\}|} \label{eq:triage}
\end{align}
The triage denominator is the set of nodes that are actually degraded when a static alarm would
have fired: the set a war room would face. Exposure is \eqref{eq:exposure} integrated over the
episode, with the remediation applied at the GNN-timed step through the verify-and-re-trigger loop,
or at the threshold-timed step.

\section{Results}
\begin{center}\small\begin{tabular}{@{}p{9.2cm}r@{}}\toprule Metric & Value \\ \midrule
Fault episodes / nominal episodes & 33 / 7 \\
Detection rate: GNN / static threshold & 1.00 / 0.97 \\
Detection rate: service $z$ (calibrated $z=4.9$ / textbook $z=3$) & 0.03 / 0.58 \\
Detection rate: per-node $z$ (calibrated $z=7.9$ / textbook $z=3$) & 0.00 / 1.00 \\
False alarms on nominal: GNN / thr. / service $z$ cal. / service $z{=}3$ / node $z$ cal. / node $z{=}3$ & 0 / 0 / 0 / 1 / 0 / 7 \\
Pre-onset false alarms on fault episodes, same order & 1 / 0 / 0 / 4 / 0 / 33 \\
Detection delay after onset (steps): GNN / static threshold & 21.6 / 35.8 \\
Detection delay (steps): service $z$ cal. / $z{=}3$; node $z$ cal. / $z{=}3$ & 38.0 / 58.6; -- / 6.9 \\
Lead time of GNN vs static threshold (steps, where it fired) & 14.2 \\
Lead time of GNN vs service $z$ (calibrated / $z{=}3$) & 18.0 / 37.2 \\
Lead time of GNN vs per-node $z$ (calibrated / $z{=}3$) & -- / -14.6 \\
Root cause hit@1 / hit@3 & 1.00 / 1.00 \\
Healed within 3 verify-and-re-trigger cycles / mean cycles & 1.00 / 1.00 \\
War-room set at alarm time (nodes) / triage reduction@3 & 15.9 / 0.33 \\
Revenue-weighted exposure: GNN-timed / threshold-timed healing & 334.2 / 815.8 \\
\bottomrule\end{tabular}\end{center}

\begin{center}\small\begin{tabular}{@{}lrrrrrr@{}}\toprule
Fault type & $n$ & delay GNN & delay thr. & delay service $z{=}3$ & hit@1 & hit@3 \\ \midrule
amf signaling storm & 9 & 25.6 & 50.2 & 58.0 & 1.00 & 1.00 \\
fiber degrade & 6 & 22.2 & 35.5 & 42.0 & 1.00 & 1.00 \\
ladn dn degrade & 1 & 20.0 & 25.0 & 31.0 & 1.00 & 1.00 \\
psa overload & 3 & 20.3 & 23.3 & 15.0 & 1.00 & 1.00 \\
router congestion & 5 & 20.4 & 35.0 & 72.0 & 1.00 & 1.00 \\
sleeping cell & 2 & 15.0 & 18.0 & 92.0 & 1.00 & 1.00 \\
upf overload & 7 & 19.4 & 32.1 & 66.8 & 1.00 & 1.00 \\
\bottomrule\end{tabular}\end{center}
\noindent\textit{Twin: 93 nodes, 198 edges. Detection threshold $\theta=0.665$. Config: episodes=160, epochs=30, T=160, window=8, seed=0. Runtime 1605.0\,s on CPU.}


\section{Reading the table}
\begin{itemize}[nosep]
  \item \textbf{Lead time.} The quantity the industry calls silent-degradation lead time is the
  GNN-versus-static-threshold row. Positive is earlier. The per-node $z$-score row is the honest
  comparison with a cheap statistical detector calibrated to the same false-alarm budget.
  \item \textbf{Root cause.} hit@1 is the fraction of episodes where the single top-ranked node is
  the injected origin, across all three domains in one softmax.
  \item \textbf{Healing cycles.} Because a remediation only removes the injection at the true
  origin, a wrong top-1 costs a cycle. The healed-within-3 rate is the end-to-end success of the
  loop, not of the model alone.
  \item \textbf{Exposure.} Revenue-weighted degradation integrated over the episode. The ratio of
  the two exposure numbers is the business version of the lead time.
\end{itemize}

\section{What would change these numbers}
The twin is small and the fault model is simple; both are the first things to vary. Longer ramps
widen the GNN's advantage; stronger coupling $\beta$ narrows it because symptoms reach the service
layer sooner. More node types and more relations make hit@1 harder. Real KPI history replaces
\eqref{eq:nominal} and will be less forgiving than AR(1) noise. None of these change the equations.
